\subsection{伴随}

在本节中， K 是域 R 或 C 之一， E 与 F 表示 K 上的希尔伯特空间。

设 \(u\) 是从 E 到 F 的稠定算子。用 \(D\) 表示 \(u\) 的定义域。对 \(y \in \mathrm{F}\) ，令 \(\lambda_y\) 为 D 上由 \(\lambda_y(x) = \langle y \mid u(x) \rangle\) （对所有 \(x \in \mathrm{D}\) ）所给的线性形式。我们用 \(\mathrm{D}^*\) 表示满足下述条件的向量 \(y \in \mathrm{F}\) 所成的集： \(\lambda_y\) 在 D 上连续。它是 F 的一个向量子空间。设 \(y \in \mathrm{D}^*\) ；由于 D 在 E 中稠密，线性形式 \(\lambda_y\) 唯一地扩张为 E 上的连续线性形式，我们仍把它记作 \(\lambda_y\) 。依照 EVT, V, p. 15，定理 3，存在唯一的元素 \(u^*(y)\) 属于 E ，使得 \(\lambda_y(x) = \langle u^*(y) \mid x \rangle\) 对所有 \(x \in \mathrm{E}\) 成立。映射 \(y \mapsto u^*(y)\) 是从 \(\mathrm{D}^*\) 到 E 中的线性映射。

定义 6. --- 以 \(D^*\) 为定义域、由 \(y \mapsto u^*(y)\) 所定义的从 F 到 E 的部分算子称为 \(u\) 的伴随，记作 \(u^*\) 。

注. --- 于是， \(y \in D^{*}\) 当且仅当存在 \(c \in R_{+}\) 使得 \(|\langle y | u(x) \rangle| \leqslant c \|x\|\) 对所有 \(x \in D\) 成立。这时元素 \(u^{*}(y)\) 由下述关系式刻画：

\[
\langle y \mid u (x) \rangle = \langle u ^ {*} (y) \mid x \rangle\tag{1}
\]

对所有 \(x \in \mathrm{D}\) 成立。这时我们有 \(|\langle y | u(x) \rangle| \leqslant \| u^{*}(y) \| \| x \|\) 对所有 \(x \in \mathrm{D}\) 成立。

在 \(u\) 是从 E 到 F 中的连续线性映射的情形，按前一定义意义的伴随与 EVT, V, p. 38，定义 1 中定义的伴随重合，因为这时 D* 等于 F 。

设 \(v \in \mathcal{P}(\mathrm{E};\mathrm{F})\) 满足 \(u \subset v\) 。这时我们有 \(v^{*} \subset u^{*}\) 。

设 \(v \in \mathcal{P}(\mathrm{E};\mathrm{F})\) 满足 \(u + v\) 稠定。这时部分算子 \(v\) 便是稠定的，且有 \(u^{*} + v^{*} \subset (u + v)^{*}\) 。一般而言等式不成立（IV, p. 347 的习题 9）。若 \(v \in \mathcal{L}(\mathrm{E};\mathrm{F})\) ，则 \(u + v\) 稠定，且 \((u + v)^{*} = u^{*} + v^{*}\) 。例如，若 \(\mathrm{F} = \mathrm{E}\) 且 \(v = \lambda 1_{\mathrm{E}}\) （其中 \(\lambda \in \mathrm{K}\) ），就是这种情形。

设 G 是 K 上的希尔伯特空间， \(v \in \mathcal{P}(\mathrm{F};\mathrm{G})\) 是稠定算子。若 \(v \circ u\) 稠定，则 \(u^{*} \circ v^{*} \subset (v \circ u)^{*}\) 。一般而言等式不成立（loc. cit.）。若 u （相应地， v ）是同构，则 \(v \circ u\) 稠定且有 \((v \circ u)^{*} = u^{*} \circ v^{*}\) 。例如，若 E = F （相应地， F = G ）且 \(u = \lambda 1_{E}\) （相应地， \(v = \lambda 1_{F}\) ），其中 \(\lambda \in K^{*}\) ，就是这种情形。

我们用 \(s\) 表示从 \(\mathrm{E} \oplus \mathrm{F}\) 到 \(\mathrm{F} \oplus \mathrm{E}\) 的希尔伯特空间等距同构，它由 \(s(x, y) = (-y, x)\) （对所有 \((x, y) \in \mathrm{E} \oplus \mathrm{F}\) ）定义。

命题 7. --- 设 u 是从 E 到 F 中的稠定部分算子。

\begin{enumerate}
\def\labelenumi{\alph{enumi})}
\item
  \(u^{*}\) 的图像等于 \(s(\Gamma_{u}^{\circ}) = s(\Gamma_{u})^{\circ}\) ；
\item
  部分算子 \(u^{*}\) 是闭的；
\item
  \(u^{*}\) 的核是 \(u\) 的像的正交补。
\end{enumerate}

记 \(\mathrm{W} = s(\Gamma_u)^\circ\) 。由于线性映射 \(s\) 是酉的，我们有 \(\mathrm{W} = s(\Gamma_u^\circ)\) 。

我们有 \((y,x)\in \mathrm{W}\) 当且仅当

\[
\langle (y, x) \mid (- u (x ^ {\prime}), x ^ {\prime}) \rangle = 0
\]

对所有 \(x' \in \mathrm{dom}(u)\) 成立，亦即若

\[
\langle y \mid u (x ^ {\prime}) \rangle = \langle x \mid x ^ {\prime} \rangle
\]

（对所有 \(x' \in \text{dom}(u)\) ）。当 \(y \in \text{dom}(u^{*})\) 且 \(x = u^{*}(y)\) 时，这一性质成立（参见 p. 236 的公式 (1)）。反之，若这一条件成立，我们推出 \(|\langle y \mid u(x') \rangle| \leqslant \|x\| \|x'\|\) 对所有 \(x' \in \text{dom}(u)\) 成立，这蕴含 \(y\) 属于 \(\text{dom}(u^*)\) ；这时我们有 \(u^*(y) = x\) 。因此 \(W = \Gamma_{u^*}\) 。

算子 \(u^{*}\) 是闭的，因为空间 \(s(\Gamma_{u})^{\circ}\) 在 \(\mathrm{F} \oplus \mathrm{E}\) 中是闭的。

我们来证明断言 \(c\) )。若 \(y\) 与 \(u\) 的像正交，则 D 上的线性形式 \(\lambda_y: x \mapsto \langle y | u(x) \rangle\) 为零，于是 \(y \in \text{dom}(u^*)\) 且 \(u^*(y) = 0\) 。反之，设 \(y \in \text{dom}(u^*)\) 。这时 \(u^*(y) = 0\) 当且仅当 \(y\) 与 \(u(x)\) （对所有 \(x \in D\) ）正交（p. 236 的公式 (1)）。

命题 8. --- 设 u 是从 E 到 F 的稠定算子。则 \(u^{*}\) 稠定当且仅当 u 可闭。在这种情形， u 的闭包 \(\overline{u}\) 等于 \(u^{**}\) ，且 \(\overline{u}\) 的伴随等于 \(u^{*}\) 。

由命题 7，部分算子 \(u^*\) 是闭的。假设 \(u^*\) 的定义域 D* 在 F 中稠密。设 \(u^{**}\) 是 \(u^*\) 的伴随；它是从 E 到 F 的闭部分算子。我们来证明 \(u \subset u^{**}\) ，这将蕴含 \(u\) 可闭。设 \(x \in \text{dom}(u)\) 。按 \(u^*\) 的定义， D* 上由 \(y \mapsto \langle x | u^*(y) \rangle\) 与 \(y \mapsto \langle u(x) | y \rangle\) 所给的线性形式相等；因此我们有 \(x \in \text{dom}(u^{**})\) 且 \(u^{**}(x) = u(x)\) ，由此即得断言。

反之，设 \(u\) 可闭；我们有 \(\Gamma_{\overline{u}} = \overline{\Gamma}_u\) （IV, p. 228 的命题 1）。设 \(y \in F\) 是与 \(\mathrm{dom}(u^*)\) 正交的一个向量。 \((y,0)\) （ \(F \oplus E\) 中的元素）这时属于 \(u^*\) 的图像的正交补。而由命题 7, a) ，我们有

\[
\Gamma_ {u ^ {*}} ^ {\circ} = (s (\Gamma_ {u}) ^ {\circ}) ^ {\circ} = \overline {{s (\Gamma_ {u})}} = s (\overline {{\Gamma}} _ {u})
\]

由此得出 \((0,y)\in\Gamma_{\overline{u}}\) ，从而 \(y=\overline{u}(0)=0\) 。由于 dom \((u^{*})\) 的正交补约化为 0 ，空间 dom \((u^{*})\) 在 F 中稠密。

最后，把命题 7 应用于 \(u^{*}\) ，即得

\[
\Gamma_ {u ^ {* *}} = s ^ {- 1} (\Gamma_ {u ^ {*}} ^ {\circ}) = s ^ {- 1} (s (\Gamma_ {u} ^ {\circ \circ})) = \overline {{\Gamma}} _ {u},
\]

因而 \(u^{**} = \overline{u}\) ，进而 \(\overline{u}^{*} = (u^{*})^{**} = \overline{u^{*}} = u^{*}\) ，因为 \(u^{*}\) 是闭的。

推论. --- 若 u 是从 E 到 F 中的闭稠定部分算子，则 \(u^{*}\) 稠定，且有 \(u^{**} = u\) 。

定义 7. --- 设 u 是 E 上的部分算子。若 u 稠定且 \(u^{*}\) 是 u 的扩张，就称 u 是对称的。若 u 稠定且 \(u^{*} = u\) ，就称 u 是自伴的。

称 u 是本质自伴的，如果它可闭且其闭包 \(\overline{u}\) 是自伴的。

称 u 是下有界部分算子，如果 u 是对称的，且存在实数 c 使得 \(\langle x | u(x) \rangle \geqslant c \|x\|^{2}\) 对属于 u 的定义域的所有 x 成立。这时称 c 是 u 的一个下界。若 c = 0 ，也称 u 是正部分算子。

我们用 \(\mathcal{A}(\mathrm{E})\) 表示 E 上自伴部分算子所成的集。

注. --- 1) 算子 \(u\) 在 E 上稠定且对称，当且仅当下式成立：

\[
\langle x \mid u (y) \rangle = \langle u (x) \mid y \rangle\tag{2}
\]

（对所有 \((x,y)\in\mathrm{dom}(u)^{2}\) ）。这一公式表明 u 的定义域包含在 \(u^{*}\) 的定义域中，进而 \(u^{*}\) 与 u 在 u 的定义域上重合。特别地，由此得出 \(\langle x|u(x)\rangle\in\mathbf{R}\) 对所有 \(x\in\mathrm{dom}(u)\) 成立。

正如在多个例子中将看到的那样，公式 (2) 常常可以通过直接计算加以验证。另一方面，伴随的定义域的精确确定——唯有它才能判别一个对称算子是否自伴——可能是非常精细的。

\begin{enumerate}
\def\labelenumi{\arabic{enumi})}
\setcounter{enumi}{1}
\item
  自伴部分算子 u 是本质自伴的（参见命题 8）。
\item
  设 \(u\) 是 E 上的对称部分算子。算子 \(u\) 可闭（命题 7, b)）。它满足 \(\text{dom}(u) \subset \text{dom}(u^*)\) ，且 \(u\) 自伴当且仅当 \(\text{dom}(u) = \text{dom}(u^*)\) 。此外， \(\overline{u}\) 的闭包 \(u\) 是对称的，因为 \(\overline{u} \subset u^* = \overline{u}^*\) （命题 8）。
\item
  设 K = C 。设 \(u \in \mathcal{P}(\mathrm{E};\mathrm{E})\) 是稠定部分算子。条件 \(\langle x | u(x) \rangle \in \mathbf{R}\) （对所有 \(x \in \operatorname{dom}(u)\) ）蕴含 u 是对称的（EVT, V, p. 2，注）；特别地，若 \(\langle x | u(x) \rangle \in \mathbf{R}_{+}\) 对所有 \(x \in \operatorname{dom}(u)\) 成立，则 u 是正的。
\item
  设 u 与 v 是 E 上的对称部分算子。若 u 自伴且 \(u \subset v\) ，则 \(v \subset v^{*} \subset u^{*} = u\) ，因此 u = v 。
\item
  本质自伴部分算子 \(u\) 是对称的，因为 \(u \subset \overline{u}\) 蕴含 \(\overline{u} = \overline{u}^{*} \subset u^{*}\) ，从而 \(u \subset u^{*}\) 。
\item
  设 u 与 v 是 E 上的对称部分算子。若 \(u+v\) 稠定（例如 u 或 v 属于 \(\mathcal{L}(\mathrm{E})\) ），则 \(u+v\) 是对称的。一般而言，部分算子 \(u + v\) 不是自伴的，即使 u 与 v 都自伴（IV, p. 347 的习题 9）。
\item
  设 \(u\) 是 E 上的对称部分算子。实数 \(c\) 是 \(u\) 的下界，当且仅当算子 \(u - c \cdot 1_{\mathrm{E}}\) 是正的。
\end{enumerate}

引理 3. --- 假设 K = R 。设 u 是从 E 到 F 中的稠定部分算子。

\begin{enumerate}
\def\labelenumi{\alph{enumi})}
\item
\(u_{(\mathbf C)}\) 的伴随是 \((u^*)_{(\mathbf C)}\) ；
\item
  设 E = F ；部分算子 u 是对称的（相应地，自伴的），当且仅当部分算子 \(u_{(\mathbf{C})}\) 是对称的（相应地，自伴的）。
\end{enumerate}

我们来证明 \(a\) )。设 \(y \in \mathrm{F}_{(\mathbf{C})}\) ，并把它写成 \(y = y_1 + iy_2\) ，其中 \(y_1, y_2 \in \mathrm{F}\) 。对所有 \((x_1, x_2) \in \mathrm{E} \times \mathrm{E}\) ，我们有

\[
\begin{array}{c} \langle u _ {(\mathbf {C})} (x _ {1} + i x _ {2})   |   y \rangle = \langle u (x _ {1})   |   y _ {1} \rangle + i \langle u (x _ {1})   |   y _ {2} \rangle \\ - i \langle u (x _ {2})   |   y _ {1} \rangle + \langle u (x _ {2})   |   y _ {2} \rangle . \end{array}
\]

若 \(y \in \text{dom}(u^*)(\mathbf{C})\) ，我们推出 \(y \in \text{dom}((u_{(\mathbf{C})})^*)\) 且 \(u_{(\mathbf{C})}^*(y) = (u_{(\mathbf{C})})^*(y)\) ，从而 \(u_{(\mathbf{C})}^* \subset (u_{(\mathbf{C})})^*\) 。

反之，设 \(y\) 属于 \(\text{dom}((u_{(\mathbf{C})})^*)\) 。在上式中取 \(x_2 = 0\) （相应地， \(x_1 = 0\) ），即可验证 \(y_1 \in \text{dom}(u^*)\) （相应地， \(y_2 \in \text{dom}(u^*)\) ），从而 \(y \in \text{dom}(u^*)(\mathbf{C})\) 。

断言 \(a\) ) 蕴含：若 \(u_{(\mathbf{C})}\) 是对称的（相应地，自伴的），则 \(u\) 也是对称的（相应地，自伴的）。

反之，设 \(u_{(\mathbf{C})}\) 是对称的。关系式 \(\langle u(x)|y\rangle = \langle x|u(y)\rangle\) （对所有 \((x,y) \in \mathrm{dom}(u_{(\mathbf{C})}) \times \mathrm{dom}(u_{(\mathbf{C})})\) ）蕴含：在取 x 与 y 属于 \(\mathrm{dom}(u)\) （ \(\mathrm{dom}(u_{(\mathbf{C})})\) 的子空间）时， u 是对称的。若 \(u_{(\mathbf{C})}\) 自伴，则前述论证表明 u 是对称的；由于 \(\mathrm{dom}(u^{*}) = \mathrm{dom}(u_{(\mathbf{C})}^{*}) \cap \mathrm{F}\) ，断言 a) 蕴含 \(\mathrm{dom}(u^{*}) = \mathrm{dom}(u_{(\mathbf{C})}) \cap \mathrm{F} = \mathrm{dom}(u)\) ，因此 u 是自伴的。

